\documentclass[12pt,a4paper]{article}

\usepackage[utf8]{inputenc}
\usepackage[T1]{fontenc}
\usepackage[french]{babel}

\usepackage{amsmath, amssymb}

\usepackage{graphicx}
\usepackage{caption}

\usepackage{siunitx} % Pour les unités

\usepackage{geometry}
\geometry{top=2cm, bottom=2.5cm, left=2cm, right=1.5cm}
\setlength{\parindent}{0pt}


\usepackage{enumitem}

\usepackage{hyperref}

\usepackage{xcolor}

\usepackage{multicol}
\setlength{\columnseprule}{1pt}
\def\columnseprulecolor{\color{gray}}


\usepackage{import}
\usepackage{xifthen}
\usepackage{pdfpages}
\usepackage{transparent}
\usepackage{float}

\newcommand{\incfig}[1]{%
	\def\svgwidth{\columnwidth}
	\import{./figures/}{#1.pdf_tex}
}



\title{Plasticité et endommagement}
\author{}
\date{\today}

\begin{document}
	
	
	\maketitle
	
	\tableofcontents
	
	\newpage
	
	\part{Chapitre 1}
	
	\section{Définition}
	
	On dit qu'un matériaux et vierge de tout endommagement si il est démunie de toute cavité et de fissures a l'échelle microscopique,
	si le matériau 'a un coportement 'a la déformation identique a celui élaboré dans les meilleurs conditions de laboratoir.
	
	\section{Aspect phénoménologique}
	
	phénoménologique : Concerne l'étude des phénomènes qui se base sur l'observation des choses telles qu'elles nous apparaissent
	
	\vspace{\baselineskip}
	
	Aucune variable mécanique d'endommagement ne peut distinguer entre un matériaux forttement endommager et d'un autre légérement ou non endommager.
	
	\section{Objectif de la théorie d'endommagement}
	
	Est de crée un variable d'endommagement noté \textbf{"D"} qui perment de quantifier l'état de la détérioration.
	
	\vspace{\baselineskip}
	
	Les mesures a l'échelle de la microstructure \big(mesure de microfissure et de densité(cavité)\big) conduisant a des modules microscopique qu'on peut intégrer a l'échelle macroscopique a l'aide des technique mathématiques d'homogéinisation.
	
	\vspace{\baselineskip}
	
	Un autre type d'évaluation d'un endommagement d'un matériau qui est lié a la durée de vie résiduelle (ce qui lui reste a vivre)
	mais ce concepte ne conduit pas directement a la notion de loi de comportement de l'endommagement.
	
	\vspace{\baselineskip}
	
	Les mesures physiques globales (résistance...) nécessite la définition d'un modéle globale pour les convertire au caractéres mécanique
	
	exemple :
	
	
	
	\begin{figure}[H]
		\centering
		\includegraphics[width=0.4\linewidth]{"figures/chapitre 1/courbe-exemple"}
		\caption{exemple sur des courbes}
		\label{fig:courbe-exemple}
	\end{figure}
	
	
	\section{Variable d'endommagement \textbf{"$D$"}}
	
	La notion de contrainte effective ou bien la section effective a etait élaboré et introduite par ROBOTNOV, elle permet de déterminer la variable d'endommagement "$D$".
	
	Soit l'aire d'une section d'un élément repérer par sa normale $ \vec{n}$ \textbf{\big( La normale est perpendiculaire a la section de la coupe \big)}
	
	Dans cette section les fissures et les vides qui constituent les dommages laisse des traces de divérse forme.
	
	\vspace{\baselineskip}
	
	Soit \textbf{$\overline{S}$} est la section effective qui résiste a l'effort. (sans vides ou fissurer (sans dommages))
	
	et \textbf{S} la section totale (avec des dommages)
	
	\[ \overline{S} \leq S \]
	
	$S_{p}$ La section totale des vides.
	
	\[ S_{p} = S-\overline{S} \]
	
	Par définition :
	
	\[ D_{\vec{n}} = \dfrac{S_{p}}{S} \quad \text{avec : } 0 \leq D_{\vec{n}} \leq 1 \]
	
	et $D_{\vec{n}}$ qui est la proportion des vides (pourcentage).
	
	\[ \text{Quand } D_{\vec{n}} \overset{\text{approche ou egale }}{\rightarrow} 1 \Rightarrow S_{p} \rightarrow S \Rightarrow \text{ Matériau trés endommager} \]
	
	\[ D_{\vec{n}} \rightarrow 0 \Rightarrow S_{p} \rightarrow 0 \Rightarrow \text{ Matériau non ou trés peu endommager} \]
	
	\[ D_{\vec{n}} \neq 0 \text{ Matériau endommager} \]
	
	En génerale quand $D_{\vec{n}} \leq 0.5$ on peut dire que la structue est récupérable ou bien non endommager.
	
	d'un point de vue physique la variable d'endommagement  est donc l'aire relative des vides coupé par la normale $\vec{n}$.
	
	\section{Hyphothése d'isopropie d'endommagement}
	
	\subsection{Définition d'isotropie}
	
	En science isotropie veut dire : qui a les mêmes propriétés dans toutes les directions.
	
	\subsection{L'hypothése}
	
	Lorsque la section des vides est la meme dans les deux sens c a d les fissures sont répartie uniformément dans les deux directions.
	
	\[ \forall \vec{n} \Rightarrow D_{\vec{n}} = D = \dfrac{S_{p}}{S}\]
	
	
	\begin{figure}[H]
		\centering
		\includegraphics[width=0.7\linewidth]{"figures/chapitre 1/endommagement-isotrope"}
		\caption{différence entre un endommagement normale et isotrope}
		\label{fig:endommagement-isotrope}
	\end{figure}
	
	
	
	
	
	
	
	
	
	
	
	
	
	
	
	
	
	
	
	
	
	
	
	
	
	
	
	
	
	
	
\end{document}